\begin{statement}{A.4.5}{Proved}{The full torsion-free cotangent assertion for irreducible normal-crossings torus fibres}
Let $g:Y\to C$ have the hypotheses of A.4.4 and let $S$ be a singular irreducible reduced full fibre with normal crossings. For every nontrivial unitary-flat line bundle $A$ on $S$,
\[
H^0(S,\widetilde\Omega_S^1\otimes A)=0,
\qquad
H^0(S,\Omega_Y^1|_S\otimes A)
 \simeq H^0(\widetilde S,\eta^*A).
\tag{A.17}
\]
The latter group has dimension one when $\eta^*A$ is holomorphically trivial, and zero otherwise. The normalization is exactly one of the following: a diagonal Hopf surface with two disjoint elliptic curves identified; a ruled surface over an elliptic curve with two disjoint sections identified by translation; or a rational surface with an anticanonical hexagon as in Schr\"oer--Siebert Proposition~2.6.
\end{statement}
\begin{proof}
We check the hypotheses of the classification, including its deformation sheaf. A.4.4 gives $\omega_S\simeq\OO_S$. Dualizing the locally free conormal resolution of $\Omega_S^1$ yields
\[
 \mathcal T^1_S=\mathcal Ext^1_{\OO_S}(\Omega_S^1,\OO_S)
 \simeq\OO_S/\mathcal J=\OO_{S_{\rm sing,red}}.
\]
Here the conormal is trivial because $S$ is the full reduced fibre; the local identities $\mathcal J=(x,y)$ at a double point and $\mathcal J=(xy,xz,yz)$ at a triple point show that this is the reduced singular locus. Hence $S$ is d-semistable in the exact sense of \cite[Definition~2.1]{SS2002}. That definition imposes irreducibility, normal crossings, d-semistability and trivial dualizing sheaf; it does not require a prior smoothing to Kodaira surfaces. Propositions~2.4, 2.5 and 2.6 there give the three normalization cases. In the ruled case their Proposition~4.1 says precisely that trivial $\omega_S$ forces the section identification to be a translation. The comparison with their notation is $X_{\mathrm{SS}}=S$ and $S_{\mathrm{SS}}=\widetilde S$.

In the rational case $\eta^*A$ is trivial: its first Chern class is torsion because the line is unitary-flat, whereas the integral second cohomology of a smooth rational surface is torsion-free and its Picard group injects into that group by the exponential sequence. Such a surface has no holomorphic one-forms, by birational invariance as proved in A.4.1. The injection $\widetilde\Omega_S^1\to\eta_*\Omega^1_{\widetilde S}$ therefore proves the first vanishing.

In the elliptic ruled case the proof of A.4.3 applies to a decomposable ruled surface of any degree: the two disjoint sections express it as $\PP_E(\OO_E\oplus P)$, with $P\in\Pic(E)$. The integral second cohomology of this ruled surface is torsion-free, by the projective-bundle formula. A unitary-flat line has torsion first Chern class, so its pullback is topologically trivial and comes from $\Pic^0(E)$. No degree assumption on $P$ occurs in the cotangent calculation, and A.4.3 proves the required vanishing for every nontrivial $A$.

In the Hopf case, pull back to a finite primary diagonal Hopf cover. Its universal covering is $\C^2\setminus\{0\}$ with contraction $(z_1,z_2)\mapsto(a_1z_1,a_2z_2)$, where $|a_i|<1$. A holomorphic one-form valued in a unitary-flat line lifts to $f_1(z)dz_1+f_2(z)dz_2$, with each $f_i$ holomorphic on the punctured space and hence on $\C^2$ by Hartogs' theorem. For a contraction character $\rho$ of modulus one, equivariance says that the coefficient of $z_1^{k_1}z_2^{k_2}dz_i$ can be nonzero only when
\[
 a_1^{k_1}a_2^{k_2}a_i=\rho
\]
(or its inverse, according to the convention for the associated line). The modulus on the left is strictly less than one and the character has modulus one. Thus every coefficient is zero. Finite etale pullback is injective on sections, so the twisted one-form group on the Hopf normalization also vanishes, and again the normalization injection proves the first assertion on $S$.

Finally A.1.2 and the first vanishing identify the ambient group with $H^0(\widetilde S,\eta^*A)$. A section of a unitary-flat line on a compact connected complex manifold is either zero or parallel and nowhere zero: its squared norm is locally a sum of squares of holomorphic functions, hence plurisubharmonic, and the maximum principle makes it constant. It therefore trivializes the line, and in the trivial case the section space is $\C$. This proves the dimension statement.
\end{proof}


% A.4.6 from part_A.tex

